\section{Matching Algorithm Pseudocode}
\label{sec:matching-pseudocode}

\begin{figure}[h]
{\footnotesize
\begin{verbatim}
MATCH(findings F, ground_truth G):
  matched_ids <- {}
  results <- []
  for each f in F:
    candidates <- []
    for each g in G:
      if g.id in matched_ids: continue
      if FILE_EQ(f, g) and CWE_IN(f, g)
         and LINE_OK(f, g):
        candidates.append(g)
    if candidates not empty:
      sort candidates: vuln=true first
      best <- candidates[0]
      matched_ids.add(best.id)
      if best.is_vulnerable:
        results.append(TP, best.id, f)
      else:
        results.append(FP, best.id, f)
    else:
      results.append(FP, null, f)
  for each g in G:
    if g.id not in matched_ids:
      if g.is_vulnerable:
        results.append(FN, g.id, null)
      else:
        results.append(TN, g.id, null)
  return results
\end{verbatim}
}
\caption{Pseudocode for the three-field matching algorithm. \texttt{FILE\_EQ} checks normalized path equality, \texttt{CWE\_IN} checks $\mathit{cwe}_f \in \mathit{cwes}_g$, and \texttt{LINE\_OK} checks line tolerance $\delta=10$.}
\label{alg:matching}
\end{figure}

\section{Ethics Statement}

All repositories evaluated in this benchmark are intentionally vulnerable
applications created explicitly for security education and training (e.g.\
PyGoat, DVPWA, VAmPI).  No responsible disclosure obligations arise from their
inclusion: the vulnerabilities they contain are by design, publicly documented,
and carry no risk of enabling attacks on production systems.  We publish no
zero-day vulnerability information and no novel exploitation techniques;
our ground-truth labels describe vulnerability \emph{classes} (CWEs) rather
than working exploit code.

Scanner weaknesses identified in our evaluation are disclosed in aggregate to
advance the state of the field.  We believe transparency about tool limitations
serves defenders more than it serves attackers, who can already assess scanner
efficacy empirically.

We acknowledge a potential conflict of interest: this benchmark is developed by
Kolega.Dev, a company that provides security tooling.  We mitigate this bias by
open-sourcing all benchmark artifacts (ground truth, scoring code, raw
scanner outputs, and evaluation harness) so that any party can audit our
methodology, re-run every evaluation, and verify reported scores independently.
Scanner results for Kolega.Dev-affiliated tools are held to the same evaluation
pipeline as all other entrants, with no post-hoc adjustment.

\section{Data Availability}

All artifacts described in this paper are publicly available under open licenses.

\begin{itemize}
  \item \textbf{Benchmark code} (MIT License) ---
    Scoring pipeline, parser library, and dashboard generator.

  \item \textbf{Ground truth} ---
    Hand-labeled vulnerability annotations for all 26 repositories.

  \item \textbf{Scan results} ---
    Raw output from all 15 evaluated scanners across every repository.

  \item \textbf{LLM evaluation harness} ---
    Prompt templates, agentic runner scripts, and per-run cost tracker.

  \item \textbf{Live dashboard} ---
    Interactive results explorer at \url{https://realvuln.kolega.dev}.

  \item \textbf{Benchmark manifest} ---
    Version-locked reproducibility file ensuring exact score reproduction.
\end{itemize}

\noindent All of the above are accessible at:
\url{https://github.com/kolega-ai/RealVulnBenchmark}
